\documentclass[10pt,a4paper]{amsart}
\usepackage{amsmath,amssymb,amsthm,graphicx}
% Retained preamble for the Stacks Project chapter snapshot in this directory.
% The chapter's own preamble.tex is a copy of the shared Stacks Project preamble; the block
% below is sliced MECHANICALLY, line for line, from it. Nothing was added, removed or
% redefined; the wrapper class, the xr-hyper/hyperref wiring and the externaldocument list
% are omitted because the assembler supplies the document class.
%
% lines taken from preamble.tex: 3, 5, 9, 11, 12, 156, 157, 158, ... (51 lines total)
%
% package coverage restored from the chapter's own preamble (2026-10-13): xy, xy-2cell

\usepackage[all]{xy}
\xyoption{2cell}
\UseAllTwocells

\newenvironment{reference}{\comment}{\endcomment}
\newenvironment{slogan}{\comment}{\endcomment}
\newenvironment{history}{\comment}{\endcomment}
\theoremstyle{plain}
\newtheorem{theorem}[subsection]{Theorem}
\newtheorem{proposition}[subsection]{Proposition}
\newtheorem{lemma}[subsection]{Lemma}
\theoremstyle{definition}
\newtheorem{definition}[subsection]{Definition}
\newtheorem{example}[subsection]{Example}
\newtheorem{exercise}[subsection]{Exercise}
\newtheorem{situation}[subsection]{Situation}
\theoremstyle{remark}
\newtheorem{remark}[subsection]{Remark}
\newtheorem{remarks}[subsection]{Remarks}
\numberwithin{equation}{subsection}
\def\lim{\mathop{\mathrm{lim}}\nolimits}
\def\colim{\mathop{\mathrm{colim}}\nolimits}
\def\Spec{\mathop{\mathrm{Spec}}}
\def\Hom{\mathop{\mathrm{Hom}}\nolimits}
\def\Ext{\mathop{\mathrm{Ext}}\nolimits}
\def\SheafHom{\mathop{\mathcal{H}\!\mathit{om}}\nolimits}
\def\SheafExt{\mathop{\mathcal{E}\!\mathit{xt}}\nolimits}
\def\Sch{\mathit{Sch}}
\def\Mor{\mathop{\mathrm{Mor}}\nolimits}
\def\Ob{\mathop{\mathrm{Ob}}\nolimits}
\def\Sh{\mathop{\mathit{Sh}}\nolimits}
\def\NL{\mathop{N\!L}\nolimits}
\def\CH{\mathop{\mathrm{CH}}\nolimits}
\def\proetale{{pro\text{-}\acute{e}tale}}
\def\etale{{\acute{e}tale}}
\def\QCoh{\mathit{QCoh}}
\def\Ker{\mathop{\mathrm{Ker}}}
\def\Im{\mathop{\mathrm{Im}}}
\def\Coker{\mathop{\mathrm{Coker}}}
\def\Coim{\mathop{\mathrm{Coim}}}
\DeclareMathSymbol{\boxtimes}{\mathbin}{AMSa}{"02}
\def\QCohstack{\mathcal{QC}\!\mathit{oh}}
\def\Cohstack{\mathcal{C}\!\mathit{oh}}
\def\Spacesstack{\mathcal{S}\!\mathit{paces}}
\def\Quotfunctor{\mathrm{Quot}}
\def\Hilbfunctor{\mathrm{Hilb}}
\def\Curvesstack{\mathcal{C}\!\mathit{urves}}
\def\Polarizedstack{\mathcal{P}\!\mathit{olarized}}
\def\Complexesstack{\mathcal{C}\!\mathit{omplexes}}
\def\Pic{\mathop{\mathrm{Pic}}\nolimits}
\def\Picardstack{\mathcal{P}\!\mathit{ic}}
\def\Picardfunctor{\mathrm{Pic}}
\def\Deformationcategory{\mathcal{D}\!\mathit{ef}}
\usepackage{cleveref}
\crefname{originallemma}{original lemma}{original lemmas}
\Crefname{originallemma}{Original lemma}{Original lemmas}
\makeatletter
\newlabel{definition-versal}{{07XF}{source}{}{sourcepointer-1}{}}
\newlabel{definition-versal@cref}{{[originallemma][0][]07XF}{[1][0][]source}{}{}{}}
\makeatother
\crefname{originallemma}{original lemma}{original lemmas}
\Crefname{originallemma}{Original lemma}{Original lemmas}
\makeatletter
\newlabel{equation-approximation}{{07X6}{source}{}{sourcepointer-2}{}}
\newlabel{equation-approximation@cref}{{[originallemma][0][]07X6}{[1][0][]source}{}{}{}}
\makeatother
\crefname{originallemma}{original lemma}{original lemmas}
\Crefname{originallemma}{Original lemma}{Original lemmas}
\makeatletter
\newlabel{formal-defos-definition-completion-CLambda}{{06GW}{source}{}{sourcepointer-3}{}}
\newlabel{formal-defos-definition-completion-CLambda@cref}{{[originallemma][0][]06GW}{[1][0][]source}{}{}{}}
\makeatother
\crefname{originallemma}{original lemma}{original lemmas}
\Crefname{originallemma}{Original lemma}{Original lemmas}
\makeatletter
\newlabel{formal-defos-remark-formal-objects-yoneda}{{06HC}{source}{}{sourcepointer-4}{}}
\newlabel{formal-defos-remark-formal-objects-yoneda@cref}{{[originallemma][0][]06HC}{[1][0][]source}{}{}{}}
\makeatother
\title{Versality of x at u\_0 matches versality of the formal object}
\author{The Stacks Project Authors}
\date{stacks-project:artin.tex:local-snapshot-20261012; GNU FDL 1.2 or later}
\begin{document}
\maketitle
\noindent\textit{Editorial note.} Statement span is lines 1902-1914 of sources/stacks-artin/pinned-original-source.tex; proof span is lines 1916-1949, the complete author proof (32 body lines).\par\medskip
\noindent\textit{Editorial note.} Navigation stubs: 4 cited-but-not-included labels are supplied as non-load-bearing pointers with their pinned tags where the snapshot tag index has them.\par\medskip
\noindent\textit{Editorial note.} extra\_packages is empty on purpose: the retained preamble already loads xy with options and a second declaration clashes in the assembler wrapper.\par\medskip
\noindent\textit{Editorial note.} Text is retained verbatim, including the author's own typos and placeholders.\par\medskip
\noindent\textit{Editorial note.} Queue index 33 of queue\_artin.json (pb=1251, stmt 1902-1914, proof 1916-1949).\par\medskip
\noindent\textit{Editorial note.} stacks-artin wave ar01 (2026-10-14): fresh adjudication accept per lane-s3/artin\_judgments.md; pool per Chief 10:5x direct dispatch (idle-first over S12's 01:5x note, SMC precedent); artin.tex 99 rows/67 with proof judged 46 accept / 21 reject / 5 AGP-overlap avoidance rows; patch-sd-envs contract spans scanned.\par\medskip
\noindent\textit{Allowed standard prerequisites (editorial).}\begin{itemize}
\item Explicit prerequisite boundary. The cited inputs are prerequisites; the argument in the retained span is not replaced by any cited result.
\end{itemize}\medskip
\ifdefined\hypertarget\hypertarget{sourcepointer-1}{}\fi
\noindent\textit{Source reference.} Original lemma 07XF: sources/stacks-artin/pinned-original-source.tex, source lines 1887--1896. Cited by the retained statement or proof but not part of the collected spans; supplied only as a navigation pointer. The argument does not depend on its content.\par\medskip
\ifdefined\hypertarget\hypertarget{sourcepointer-2}{}\fi
\noindent\textit{Source reference.} Original lemma 07X6: sources/stacks-artin/pinned-original-source.tex, source lines 1289--1295. Cited by the retained statement or proof but not part of the collected spans; supplied only as a navigation pointer. The argument does not depend on its content.\par\medskip
\ifdefined\hypertarget\hypertarget{sourcepointer-3}{}\fi
\noindent\textit{Source reference.} Original lemma 06GW: sources/stacks-artin/pinned-original-source.tex, source lines 1949--1949. Cited by the retained statement or proof but not part of the collected spans; supplied only as a navigation pointer. The argument does not depend on its content.\par\medskip
\ifdefined\hypertarget\hypertarget{sourcepointer-4}{}\fi
\noindent\textit{Source reference.} Original lemma 06HC: sources/stacks-artin/pinned-original-source.tex, source lines 1949--1949. Cited by the retained statement or proof but not part of the collected spans; supplied only as a navigation pointer. The argument does not depend on its content.\par\medskip
\begin{lemma}
\label{lemma-versality-matches}
With notation as in Definition \ref{definition-versal}.
Let $R = \mathcal{O}_{U, u_0}^\wedge$.
Let $\xi$ be the formal object of $\mathcal{X}$
over $R$ associated to $x|_{\Spec(R)}$, see (\ref{equation-approximation}).
Then
$$
x\text{ is versal at }u_0
\Leftrightarrow
\xi\text{ is versal}
$$
\end{lemma}

\begin{proof}
Observe that $\mathcal{O}_{U, u_0}$ is a Noetherian local $S$-algebra
with residue field $k$. Hence $R = \mathcal{O}_{U, u_0}^\wedge$ is an object of
$\mathcal{C}_\Lambda^\wedge$, see Formal Deformation Theory, Definition
\ref{formal-defos-definition-completion-CLambda}.
Recall that $\xi$ is versal if
$\underline{\xi} : \underline{R}|_{\mathcal{C}_\Lambda} \to
\mathcal{F}_{\mathcal{X}, k, x_0}$
is smooth and $x$ is versal at $u_0$ if
$\hat x : \mathcal{F}_{(\Sch/U)_{fppf}, k, u_0}
\to \mathcal{F}_{\mathcal{X}, k, x_0}$ is smooth.
There is an identification of predeformation categories
$$
\underline{R}|_{\mathcal{C}_\Lambda}
=
\mathcal{F}_{(\Sch/U)_{fppf}, k, u_0},
$$
see Formal Deformation Theory, Remark
\ref{formal-defos-remark-formal-objects-yoneda} for notation.
Namely, given an Artinian local $S$-algebra $A$ with residue field
identified with $k$ we have
$$
\Mor_{\mathcal{C}_\Lambda^\wedge}(R, A) =
\{\varphi \in \Mor_S(\Spec(A), U) \mid \varphi|_{\Spec(k)} = u_0\}
$$
Unwinding the definitions the reader verifies that the resulting map
$$
\underline{R}|_{\mathcal{C}_\Lambda} =
\mathcal{F}_{(\Sch/U)_{fppf}, k, u_0}
\xrightarrow{\hat x}
\mathcal{F}_{\mathcal{X}, k, x_0},
$$
is equal to $\underline{\xi}$ and we see that the lemma is true.
\end{proof}

\end{document}
